\documentclass[12pt,a4paper]{article}

\usepackage[T1]{fontenc}
\usepackage[utf8]{inputenc}
\usepackage{amsmath,amssymb,amsthm}
\usepackage{mathtools}
\usepackage{geometry}
\usepackage{hyperref}
\usepackage{booktabs}
\usepackage{array}
\usepackage{xcolor}
\usepackage{tikz-cd}
\usepackage{natbib}

\geometry{margin=2.5cm}

\hypersetup{
  colorlinks=true,
  linkcolor=blue!70!black,
  citecolor=blue!70!black,
  urlcolor=blue!70!black
}

\newtheorem{theorem}{Theorem}[section]
\newtheorem{lemma}[theorem]{Lemma}
\newtheorem{corollary}[theorem]{Corollary}
\newtheorem{definition}[theorem]{Definition}
\newtheorem{conjecture}[theorem]{Conjecture}
\newtheorem{observation}[theorem]{Observation}
\theoremstyle{remark}
\newtheorem{remark}[theorem]{Remark}
\newtheorem{openproblem}[theorem]{Open Problem}

\newcommand{\PDL}{\textsc{pdl}}
\newcommand{\Coh}{\mathrm{Coh}}
\newcommand{\Ztwo}{\mathbb{Z}_{2}^{3}}
\newcommand{\Qphi}{\mathbb{Q}(\sqrt{5})}
\newcommand{\ksurf}{\kappa_{\mathrm{surf}}}
\newcommand{\miso}{\Delta m_{\mathrm{iso}}}
\newcommand{\Ntot}{N_{\mathrm{tot}}}
\newcommand{\thetaW}{\theta_{W}}
\newcommand{\sinthW}{\sin^{2}\!\theta_{W}}
\newcommand{\cosW}{\cos\theta_{W}}
\newcommand{\Dalep}{\Delta\alpha_{\mathrm{lep}}}
\newcommand{\Dahad}{\Delta\alpha_{\mathrm{had}}}
\newcommand{\MZ}{M_{Z}}
\newcommand{\MW}{M_{W}}
\newcommand{\MH}{M_{H}}
\newcommand{\vev}{v}
\newcommand{\alphaPDL}{\alpha_{\mathrm{PDL}}}
\newcommand{\ppm}{\,\mathrm{ppm}}
\newcommand{\MeV}{\,\mathrm{MeV}}
\newcommand{\GeV}{\,\mathrm{GeV}}
\newcommand{\TeV}{\,\mathrm{TeV}}
\newcommand{\mproton}{m_{p}}

\pdfstringdefDisableCommands{%
  \def\PDL{PDL}%
  \def\Ztwo{Z2\^{}3}%
  \def\MZ{MZ}%
  \def\MW{MW}%
  \def\MH{MH}%
  \def\thetaW{theta\_W}%
  \def\ksurf{kappa\_surf}%
}

\title{%
  \textbf{D62: Gauge Boson Masses from Combinatorial Axioms}\\[0.4em]
  \large Derivation of \texorpdfstring{$\vev$, $\MZ$, $\MW$}{v, MZ, MW}
  and Prediction of \texorpdfstring{$\MH$}{MH}
  in the Projective Dynamic Logo Framework
}

\author{C\'edric Laubscher\\[0.3em]
  \small Independent Researcher, Switzerland\\
  \small ORCID: 0009-0004-5415-1098\\
  \small \href{https://cedriclaubscher.ch}{cedriclaubscher.ch}
}

\date{June 2026}

\begin{document}

\maketitle

\begin{abstract}
Building on the derivation of the Standard Model gauge group
$\mathrm{SU}(3)\times\mathrm{SU}(2)\times\mathrm{U}(1)$ as an unconditional
theorem of the four combinatorial axioms C1--C4 of the Projective Dynamic Logo
(\PDL) programme (D57--D61), this document addresses the question of gauge
boson masses. We provide a formal definition of a gauge boson mass as the
minimal coherence excitation cost of a transport mode in
$\Coh(K_4)\cong\Ztwo$. The electroweak vacuum expectation value is derived as
$\vev = \Ntot^{2}/(R_{e}\,k_{1})\cdot\mproton$ from the structure of the
$\mathrm{SU}(2)$ sequencing network (Conjecture~T1, $676\ppm$). The ratio
$\MZ/\MW = 1/\cosW$ is an unconditional theorem, with the deviation of
$5182\ppm$ from experiment understood as a tree-level vs.\ physical angle
effect of order $\alpha/\pi$. An analytical chain for $\MZ$ via the running
electromagnetic coupling yields $\MZ/\mproton = 97.179$ at $84.6\ppm$
($3.67\sigma$); the residual tension is traced precisely to the absence of
sub-proton hadronic closures (pions, kaons) in the current corpus, pending
resolution of OP-D59-1. The Higgs mass is predicted as
$\MH/\mproton = 133.611$ at $1.49\sigma$. The decay cascade is mapped onto
the $\Ztwo$ topology: the Higgs admits exactly four primary decompositions,
neutrinos exit the coherence network at level~2 carrying $12.2\%$ of the
Higgs rest energy, and the production cross-section admits a structural
estimate within $17871\ppm$. Five open problems are precisely formulated with
their resolution pathways.
\end{abstract}

\clearpage

\tableofcontents
\newpage

%%------------------------------------------------------------
\section{Introduction and context}
%%------------------------------------------------------------

Documents D57 through D61 of the \PDL\ corpus established that the Standard
Model gauge group $G_{\mathrm{SM}}=\mathrm{SU}(3)\times\mathrm{SU}(2)\times
\mathrm{U}(1)$ is an unconditional theorem of axioms C1--C4 on finite signed
graphs~\cite{D57,D58,D59,D60,D61}. D60 proved that $G_{\mathrm{eff}}=
S_{4}/V_{4}\cong S_{3}$ acts on $\Coh(K_4)\cong\Ztwo$ (D60), and D61
showed that the minimal covariant derivative $D_{\mu}=\partial_{\mu}-igA_{\mu}$
is the unique C4-compatible gauge transport. These results complete the gauge
sector of the \PDL\ programme without free parameters beyond $\miso=m_{d}-m_{u}$.

What remained open was the question of \emph{mass}. Deriving a gauge group
without deriving the masses of its associated bosons is structurally incomplete:
the group tells us which transport modes are allowed, but not what it costs to
excite them. In the Standard Model (SM), boson masses arise via the Higgs
mechanism~\cite{Higgs1964,Brout1964}: a scalar field acquires a vacuum
expectation value $\vev\approx246\GeV$, spontaneously breaking $\mathrm{SU}(2)
\times\mathrm{U}(1)\to\mathrm{U}(1)_{\mathrm{em}}$, so that $\MW=g\vev/2$ and
$\MZ=\MW/\cosW$. This mechanism is postulated, not derived.

The present document addresses this gap. It provides: a formal definition of
gauge boson mass in the \PDL\ formalim; a derivation of the electroweak vacuum
expectation value from C1--C4; the ratio $\MZ/\MW$ as an unconditional theorem;
an analytical chain for $\MZ$; a prediction for $\MH$; and a structural account
of Higgs production and decay in the \PDL\ language.

\paragraph{Physical picture.} In the \PDL\ paradigm, a proton--proton collision
at high energy is a \emph{coherence reorganisation}: two fermionic closures
attempt to form a common coherent state, and axiom C4 selects the
reconfiguration of minimal cost at each step. The resulting \emph{decay cascade}
is a hierarchical leakage process. At each level, the network exports its
coherence cost towards smaller-scale admissible closures; the observable
particles are the ``quotients'' of this hierarchical division, and the
undetected neutrinos and soft photons are the ``remainders'' that cannot
organise into any further admissible closure. The mass of each boson is the
threshold energy at which the corresponding reorganisation becomes accessible.

\paragraph{Epistemic convention.} Every result in this document carries an
explicit epistemic status: unconditional theorem of C1--C4, well-motivated
conjecture, structural hypothesis, or open problem. Conjectures are presented
with their numerical tension against experiment. No result is claimed beyond
what the derivation establishes.

%%------------------------------------------------------------
\section{Formal definition of gauge boson mass in \texorpdfstring{\PDL}{PDL}}
%%------------------------------------------------------------

\subsection{\texorpdfstring{The coherence network $\Coh(K_4)\cong\Ztwo$}{The coherence network}}

From D60, the coherence group of $K_4$ is $\Coh(K_4)\cong\Ztwo$, with eight
elements and seven non-trivial generators. The group $G_{\mathrm{eff}}=S_{3}$
acts on $\Ztwo$ by permuting the three bits, partitioning the seven generators
into three orbits:
\begin{align*}
\mathcal{O}_{1} &= \{(1,0,0),\,(0,1,0),\,(0,0,1)\}, &|\mathcal{O}_{1}|=3,\\
\mathcal{O}_{2} &= \{(1,1,0),\,(1,0,1),\,(0,1,1)\}, &|\mathcal{O}_{2}|=3,\\
\mathcal{O}_{3} &= \{(1,1,1)\},                       &|\mathcal{O}_{3}|=1.
\end{align*}
By D58 and D59, $\mathcal{O}_{1}$ corresponds to the gluon sector (massless,
commutes with the C1 binary pulsation), $\mathcal{O}_{2}$ to the weak bosons
($W^{\pm}$, $Z^{0}$, massive), and $\mathcal{O}_{3}$ to the Higgs scalar
(massive). The photon is the U(1) phase mode: it commutes with C1 and carries
zero coherence cost, hence zero mass.

\subsection{Definition of gauge boson mass}

In the \PDL\ corpus, the mass of a fermionic closure is its stationary
coherence cost (D32). Gauge bosons are not fermionic closures; they are
transport conditions between closures. Their mass cannot be defined as a
stationary cost. The correct definition is instead the \emph{minimal excitation
threshold}:

\begin{definition}[Gauge boson mass in \PDL]\label{def:boson-mass}
Let $B\in\Ztwo$ be a non-trivial generator. The \emph{mass} of the
corresponding gauge boson is:
\begin{equation}
m_{B} = \inf\Bigl\{E>0\;\Big|\;\exists\,\mathcal{C}\in\Ztwo:
\nu(\mathcal{C})-\nu_{0} = \frac{E\hbar}{c^{2}R_{\mathrm{tot}}}\,,\;
\mathcal{C}\text{ activates mode }B\Bigr\},
\end{equation}
where $\nu_{0}$ is the ground-state coherence fraction of the $\mathrm{SU}(2)$
transport network and $R_{\mathrm{tot}}=11017$.
\end{definition}

The generators of $\mathcal{O}_{1}$ (gluons, photon) have $m_{B}=0$: they
correspond to symmetries of the ground state that commute with C1, so no
threshold energy is required to activate them. The generators of
$\mathcal{O}_{2}$ ($W^{\pm}$, $Z^{0}$) break C1 pulsation locally (in one
or two bit dimensions), incurring a positive threshold cost. The generator
$(1,1,1)\in\mathcal{O}_{3}$ (Higgs) breaks C1 globally in all three
dimensions, incurring the largest threshold cost.

\subsection{The decay cascade and its topological structure}

The generator $(1,1,1)\in\mathcal{O}_{3}$ admits exactly four direct
decompositions in $\Ztwo$ (verified exhaustively):
\begin{equation}\label{eq:decomp}
(1,1,1) = \underbrace{(0,0,1)+(0,1,0)+(1,0,0)}_{\text{1 decomp.\ via }\mathcal{O}_{1}}
= \underbrace{\begin{cases}
(0,0,1)+(1,1,0)\\
(0,1,0)+(1,0,1)\\
(1,0,0)+(0,1,1)
\end{cases}}_{\text{3 decomp.\ via }\mathcal{O}_{1}+\mathcal{O}_{2}}.
\end{equation}
These four paths define the four dominant primary decay channels of the Higgs
boson: $H\to ggg$ (one way), and $H\to g+W^{\pm}$ or $g+Z^{0}$ (three ways).
The cascade ordering $\mathcal{O}_{3}\to\mathcal{O}_{2}\to\mathcal{O}_{1}$
is not arbitrary: it is the unique ordering compatible with C4 (cost
minimisation). Orbits of higher weight carry higher coherence cost; C4 forces
the network to discharge towards lower-cost orbits one level at a time.

This structure has a precise analogy. In a high-energy collision, the energy
budget is partitioned hierarchically: the Higgs absorbs the coherent fraction,
the $W/Z$ bosons absorb the next level, and hadrons and leptons absorb the
remainder. At each level, C4 selects the reconfiguration of minimal cost; what
cannot be absorbed at one level is exported to the next. The observable
particles are the ``quotients'' of this hierarchical process, and the
undetected particles (soft photons, neutrinos) are the ``remainders'' that
exit the coherence network without further reorganisation.

\begin{observation}[Neutrino exit level]\label{obs:neutrinos}
Neutrinos enter the cascade at level~2, when generators of $\mathcal{O}_{2}$
decompose into pairs of $\mathcal{O}_{1}$ elements. For example,
$(0,1,1)\to(0,1,0)+(0,0,1)$. In the \PDL\ formalism (D33), the neutrino
corresponds to the neutral $\mathcal{O}_{1}$ component with
$R_{\mathrm{surf}}(\nu)\approx0$. A closure with near-zero surface budget
cannot form new coherent triangles with any other closure; it exits the
coherence network immediately upon production.

The energy fraction carried away invisibly by neutrinos from the Higgs is
estimated at $12.2\%$ of $M_{H}$, arising from $WW^{*}$ ($7.5\%$), $ZZ^{*}$
($0.5\%$), and $\tau^{+}\tau^{-}$ ($4.1\%$) channels. This fraction is
absent from the reconstructed invariant mass used by ATLAS and CMS; the
measured mass $125.2\GeV$ is thus the \emph{coherent fraction}, not the total
reorganisation cost. This observation is an unconditional structural consequence
of Definition~\ref{def:boson-mass} and D33.
\end{observation}

%%------------------------------------------------------------
\section{The electroweak vacuum expectation value}
%%------------------------------------------------------------

\subsection{\texorpdfstring{The $\mathrm{SU}(2)$ sequencing network}{The SU(2) sequencing network}}

By D55 (Lemma~D), the total number of admissible $\mathrm{SU}(2)$ sequencing
cycles in the transport network is:
\begin{equation}
\Ntot = \frac{n_{d}}{4}\cdot\frac{r_{d}-r_{u}}{R_{e}} = 7\times17=119,
\end{equation}
where $n_{d}=28$, $r_{d}=\binom{28}{2}=378$, $r_{u}=\binom{24}{2}=276$,
$R_{e}=6$. This is an unconditional theorem of C1--C4.

\begin{definition}[Coherent transports on the sequencing network]\label{def:transport}
The $\mathrm{SU}(2)$ sequencing network is the cyclic group $\mathbb{Z}_{\Ntot}$
of $\Ntot=119$ cycles. A \emph{coherent transport} is a group endomorphism
$\mathbb{Z}_{\Ntot}\to\mathbb{Z}_{\Ntot}$; there are exactly $\Ntot=119$ such
endomorphisms, including the null endomorphism (identity). By C1 (binary
pulsation), each transport simultaneously engages \emph{two} endomorphisms,
one per pulsation branch. The total capacity of the network is therefore the
number of ordered pairs of coherent transports:
\begin{equation}
\mathcal{T} = \Ntot^{2} = 119^{2} = 14161.
\end{equation}
\end{definition}

\begin{definition}[Minimal transport unit]\label{def:mtu}
The minimal coherent transport unit is the product of the elementary $K_4$
relational budget and the primary $\mathrm{SU}(2)$ leakage period:
\begin{equation}
\mathcal{U}_{\min} = R_{e}\cdot k_{1} = 6\times9 = 54,
\end{equation}
where $R_{e}=6$ (D16a) and $k_{1}=n_{d}-n_{u}+R_{e}-1=9$ (D51) are
unconditional theorems.
\end{definition}

\subsection{Derivation of the vev}

\begin{conjecture}[Electroweak vev]\label{conj:vev}
The vacuum expectation value of the electroweak symmetry breaking is:
\begin{equation}\label{eq:vev}
\vev = \frac{\Ntot^{2}}{R_{e}\,k_{1}}\cdot\mproton
= \frac{14161}{54}\cdot\mproton \approx 262.241\,\mproton.
\end{equation}
The deviation from the experimental value $\vev_{\exp}/\mproton=262.418$ is
$676\ppm$.
\end{conjecture}

\emph{Justification.} The vev is the energy threshold below which the
$\mathrm{SU}(2)$ transport network cannot sustain a fully coherent state: the
minimal energy per unit is $\mproton$ (the reference fermionic closure), the
total capacity is $\mathcal{T}=\Ntot^{2}$ ordered transport pairs, and the
unit cost of a single coherent transport is $\mathcal{U}_{\min}=R_{e}k_{1}$.
Their ratio $\Ntot^{2}/(R_{e}k_{1})$ is the number of ``unit transports''
that the network can accommodate. In an exhaustive numerical search over all
rational combinations of \PDL\ quintuplet invariants, this candidate is
isolated by a factor of 15 from the next alternative.

\begin{remark}[Correction and OP-D62-1]\label{rem:vcorr}
The residual of $676\ppm$ is real: $\vev_{\exp}$ is defined exactly via
$G_{F}$, known to $0.5\ppm$. The leading correction candidate is:
\begin{equation}\label{eq:vcorr}
\vev_{\mathrm{corr}} = \frac{\Ntot^{2}}{R_{e}\,k_{1}}
\left(1+\frac{k_{1}}{\Ntot^{2}}\right)\cdot\mproton
= \frac{14170}{54}\cdot\mproton,
\end{equation}
reducing the deviation to $41\ppm$. The factor $k_{1}/\Ntot^{2}=9/14161$ is
interpretable as the weight of the null endomorphism (identity transport) in
$\mathbb{Z}_{\Ntot}$: the null transport contributes to the network capacity
but carries no physical information, and its fractional weight
$k_{1}/\Ntot^{2}$ corrects the vev upward by the corresponding proportion.
The formal derivation of this correction from C1--C4 is OP-D62-1.
\end{remark}

%%------------------------------------------------------------
\section{Analytical derivation of \texorpdfstring{$M_Z$}{MZ}}
%%------------------------------------------------------------

\subsection{Strategy: the SM relation as a PDL chain}

In the SM, $\MZ = g_{Z}(\MZ)\,\vev/2$ where $g_{Z}=e(\MZ)/(\sin\thetaW\cos\thetaW)$
and $e(\MZ)$ is the electromagnetic coupling evaluated at the scale $\MZ$.
This relation is exact in the SM and is a consequence of the $\mathrm{SU}(2)
\times\mathrm{U}(1)$ breaking structure. In the \PDL, the relation holds with:
(i)~$\vev$ from Conjecture~\ref{conj:vev}; (ii)~$\thetaW=19\pi/119$ from D55
(unconditional theorem); (iii)~$e(\MZ)$ from the \PDL\ coupling constant
$\alphaPDL=1/137.036$ (D12) corrected for the running between $0$ and $\MZ$.

The running correction $\Delta\alpha(\MZ)=1-\alpha(0)/\alpha(\MZ)$ encodes how
the electromagnetic coupling increases from its low-energy value to the
electroweak scale. In the \PDL\ language, this increase is a \emph{granularity
effect}: at higher energies, the network engages more degrees of freedom, and
the effective coupling is larger. The correction decomposes into two structurally
distinct contributions.

\subsection{Leptonic running}

The Schwinger perturbative formula gives the leptonic contribution:
\begin{equation}\label{eq:dalep}
\Dalep = \frac{\alphaPDL}{3\pi}\sum_{\ell=e,\mu,\tau}
\!\left[\ln\frac{\MZ^{2}}{m_{\ell}^{2}}-\frac{5}{3}\right] = 0.031421,
\end{equation}
at $2447\ppm$ from the PDG value $0.031498$. This contribution is derivable
from D12 ($\alphaPDL$) and the lepton masses, both already in the corpus.
Each lepton $\ell$ contributes proportionally to $\ln(\MZ^{2}/m_{\ell}^{2})$:
the logarithm counts the number of energy ``levels'' between $m_{\ell}$ and
$\MZ$, i.e.\ the number of network scales that the transport must traverse.

\subsection{Hadronic running and the PDL lower bound}

In the SM, the hadronic contribution is the dispersive integral:
\begin{equation}\label{eq:dahad_SM}
\Dahad^{(\mathrm{SM})} = -\frac{\alpha\MZ^{2}}{3\pi}
\mathcal{P}\!\int_{4m_{\pi}^{2}}^{\infty}
\frac{R(s)}{s(s-\MZ^{2})}\,\mathrm{d}s, \qquad
R(s) = \frac{\sigma(e^{+}e^{-}\to\mathrm{hadrons})}{\sigma(e^{+}e^{-}\to\mu^{+}\mu^{-})}.
\end{equation}
The lower bound $4m_{\pi}^{2}$ is the production threshold of the lightest
charged hadronic state. In the \PDL\ framework, this threshold has a precise
structural meaning: it is the minimal energy at which a hadronic fermionic
closure can be produced. Pions and kaons, which dominate the region
$[2m_{\pi},\,m_{\rho}]$ and contribute approximately $0.003$ to
$\Dahad^{(\mathrm{SM})}$, are Goldstone modes of chiral symmetry breaking.
They are not independent C1--C4 fermionic closures in the present corpus; they
will be derived from C1--C4 in the resolution of OP-D59-1.

Consequently, the natural lower bound of the \PDL\ dispersive sum is $\mproton$:
only the proton and heavier hadrons are admissible fermionic closures in the
current corpus. With this lower bound, the hadronic contribution to the running
is the surface leakage of the reference fermionic closure, normalised by the
golden ratio:

\begin{conjecture}[Hadronic running coupling]\label{conj:dahad}
\begin{equation}\label{eq:dahad}
\Dahad = \frac{\ksurf}{\varphi} = \frac{310}{11017} \approx 0.028138,
\end{equation}
where $\ksurf=\varphi\cdot r_{\mathrm{val}}/(3\,R_{\mathrm{tot}})=\varphi\cdot310/11017$
is the surface leakage parameter of the proton (D42) and $\varphi=(1+\sqrt{5})/2$.
\end{conjecture}

\emph{Justification.} The value $\ksurf/\varphi=310/11017$ is the fraction of
the proton's surface-active relational budget divided by $\varphi$ (the
Fibonacci normalisation of the surface-to-core ratio, D05). It is the
contribution of the proton-sector fermionic closures to the electromagnetic
running when the dispersive integral is evaluated with lower bound $\mproton$
rather than $2m_{\pi}$.

\emph{Tension and its origin.} The PDG value is $\Dahad^{(\mathrm{PDG})}=
0.027661\pm0.000110$, at $4.34\,\sigma$ from $\ksurf/\varphi=0.028138$.
This tension is not an error: it is the missing contribution of pions and
kaons ($\approx0.000477$), which in the SM integral contribute from threshold
$2m_{\pi}$ but are absent from the \PDL\ lower bound. When OP-D59-1 provides
the pion and kaon contributions in \PDL\ language, the correction will bring
$\Dahad^{(\mathrm{PDL})}$ into agreement with $\Dahad^{(\mathrm{PDG})}$, and
the tension on $\MZ$ is expected to fall below $2\sigma$.

\subsection{\texorpdfstring{The $M_Z$ prediction}{MZ prediction}}

\begin{conjecture}[\texorpdfstring{$M_Z$}{MZ}]\label{conj:mz}
The $Z^{0}$ boson mass satisfies:
\begin{equation}\label{eq:mz}
\frac{\MZ}{\mproton} = \frac{g_{Z}\bigl[\Dalep+\Dahad\bigr]\cdot\vev}{2\mproton},
\end{equation}
where $\alpha_{\mathrm{eff}}=\alphaPDL/(1-\Dalep-\Dahad)$,
$g_{Z}=\sqrt{4\pi\alpha_{\mathrm{eff}}}/(\sin\thetaW\cos\thetaW)$,
$\thetaW=19\pi/119$ \emph{(D55)}, and $\vev$ is equation~\eqref{eq:vev}.
Numerically: $\MZ/\mproton=97.179$, deviation $84.6\ppm$ from experiment,
tension $3.67\sigma$.
\end{conjecture}

The self-consistency of equation~\eqref{eq:mz} was verified by fixed-point
iteration: starting from $\MZ^{(0)}=97\,\mproton$, the sequence converges in
two iterations (the rapid convergence reflects the fact that $\Dahad=
\ksurf/\varphi$ is independent of $\MZ$). The tension of $3.67\sigma$
against the experimental value $\MZ^{(\mathrm{exp})}=91187.6\MeV$ is
structural: it reflects the $4.34\sigma$ discrepancy in $\Dahad$ arising from
the absence of pion and kaon contributions. The two tensions are linked
quantitatively: $\delta(\Dahad)\times\mathrm{d}\MZ/\mathrm{d}(\Dahad)\approx
0.000477\times51.7\,\mproton\approx24.7\,\mproton$, corresponding to
$\approx11\sigma_{\MZ}$ if $\Dahad$ alone were corrected. The smaller observed
tension of $3.67\sigma$ on $\MZ$ results from the partial compensation between
the $\Dahad$ excess and the $\vev$ deficit.

%%------------------------------------------------------------
\section{The theorem \texorpdfstring{$M_Z/M_W=1/\cos\theta_W$}{MZ/MW}}
%%------------------------------------------------------------

\begin{theorem}[\texorpdfstring{$M_Z/M_W$}{MZ/MW}]\label{thm:mzmw}
As an unconditional consequence of C1--C4:
\begin{equation}\label{eq:mzmw}
\frac{\MZ}{\MW} = \frac{1}{\cosW} = \frac{1}{\cos\!\left(\tfrac{19\pi}{119}\right)}.
\end{equation}
The PDL value $1/\cosW=1.14049$ differs from the experimental ratio
$\MZ/\MW=1.13461$ by $5182\ppm$. This deviation is consistent with the known
difference between the tree-level Weinberg angle and its radiatively corrected
physical value, an effect of order $\alpha/\pi\approx2300\ppm$ in the SM.
\end{theorem}

\begin{proof}
From D57, the relation $\MW=\MZ\cosW$ is a structural consequence of
$\mathrm{SU}(2)\times\mathrm{U}(1)\to\mathrm{U}(1)_{\mathrm{em}}$ breaking,
itself an unconditional theorem of C1--C4. From D55, $\thetaW=19\pi/119$ is
an unconditional theorem. The ratio $\MZ/\MW=1/\cosW$ then holds without free
parameters. The $5182\ppm$ deviation from the experimental ratio reflects the
difference between the tree-level (PDL) and radiatively corrected (physical)
Weinberg angle. Comparing $\sinthW^{(\mathrm{PDL})}=0.23120$ against the PDG
MS-bar value $\sinthW^{(\mathrm{exp})}=0.23121\pm0.00003$ gives a tension
of $0.48\sigma$, confirming that the tree-level \PDL\ angle and the physical
angle agree to better than $1\sigma$.
\end{proof}

\begin{corollary}\label{cor:mw}
Once $\MZ$ is determined, $\MW$ is a theorem of C1--C4:
$\MW = \MZ\cdot\cos(19\pi/119)$.
\end{corollary}

%%------------------------------------------------------------
\section{Prediction of \texorpdfstring{$M_H$}{MH}}
%%------------------------------------------------------------

\begin{conjecture}[\texorpdfstring{$M_H$}{MH}]\label{conj:mh}
\begin{equation}\label{eq:mh}
\frac{\MH}{\mproton} = \frac{\varphi\,k_{2}}{\sinthW}
+ \frac{\Delta n\cdot k_{2}}{\Ntot}
= \frac{\varphi\times19}{\sin^{2}(19\pi/119)} + \frac{4\times19}{119},
\end{equation}
where $k_{2}=19$ \emph{(D51)}, $\Delta n=4$ \emph{(D47)}, $\Ntot=119$
\emph{(D55)}. Numerically: $\MH/\mproton=133.611$ at $1307\ppm$ ($1.49\sigma$).
\end{conjecture}

\emph{Structural motivation.} The generator $(1,1,1)\in\mathcal{O}_{3}$ is the
unique element of $\Ztwo$ that engages all three transport dimensions
simultaneously. Its excitation cost is therefore the largest in the network and
involves all three \PDL\ structural scales at once: $\varphi$ (surface-to-core
ratio, D05), $k_{2}=19$ ($\mathrm{SU}(2)$ sequencing cycle, D51), and
$\sinthW$ (electroweak mixing, D55). The correction term $\Delta n\cdot k_{2}/\Ntot
=76/119$ encodes the isospin asymmetry $\Delta n=4$ (D47) acting on the
$\mathrm{SU}(2)$ cycle length.

\emph{Epistemic status.} The Higgs self-coupling $\lambda_{H}=(M_{H}/v)^{2}/2
\approx0.129$ is not yet derived from C1--C4. Its derivation requires the top
quark Yukawa coupling $y_{t}=\sqrt{2}\,m_{t}/v$, hence the top quark mass from
OP-D59-1 (OP-D62-2). Until then, Conjecture~\ref{conj:mh} is a motivated
prediction, not a theorem.

%%------------------------------------------------------------
\section{Higgs production and decay: structural observations}
%%------------------------------------------------------------

\subsection{Rarity of Higgs production}

At $\sqrt{s}=8\TeV$, $\sigma(gg\to H)\approx19\,\mathrm{pb}$ while
$\sigma_{pp}\approx70\,\mathrm{mb}$, giving a ratio of approximately
$2.71\times10^{-10}$. This extraordinary rarity reflects the complexity of the
required coherence configuration. A Higgs is produced only when the network
simultaneously activates all three dimensions of $\Ztwo$, which requires:
(i)~a coherent encounter of two $K_{4}$ surfaces; (ii)~selection of the
specific $\mathcal{O}_{3}$ generator among all configurations of $\Ztwo$.

$K_{4}$ has exactly three perfect matchings. When two protons collide, the
number of jointly coherent configurations is $3\times3=9$ out of
$6\times4=24$ possible mixed triangles, giving a coherence probability of
$9/24=3/8$. Among coherent configurations, the probability of activating
$(1,1,1)$ specifically is $1/8$ (one element among eight in $\Ztwo$). The
structural estimate is:
\begin{equation}\label{eq:xsec}
\frac{\sigma(gg\to H)}{\sigma_{pp}} \sim \frac{\ksurf^{3}}{\Ntot^{2}\cdot5^{2}},
\end{equation}
where $5^{2}=(\sqrt{5})^{4}$ encodes the double encounter of two $K_4$
surfaces, each contributing a factor $\sqrt{5}$ through $\ksurf=\varphi\cdot310/
R_{\mathrm{tot}}$ and the relation $\varphi=(1+\sqrt{5})/2$. Numerically,
$\ksurf^{3}/(\Ntot^{2}\times25)=2.67\times10^{-10}$, within $17871\ppm$ of
the experimental value. The derivation of the factor $5^{2}$ from a
combinatorial matching problem on $K_4$ is OP-D62-5.

\subsection{Decay cascade: three levels}

The cascade has a precise three-level structure within $\Ztwo$:
\begin{enumerate}
\item[\textbf{Level 0.}] Production: $(1,1,1)$ is created. Lifetime
$\tau_{H}\sim\hbar/\Gamma_{H}\approx10^{-22}\,\mathrm{s}$. No stable
intermediate state exists at this orbit level.
\item[\textbf{Level 1.}] Primary decay: $(1,1,1)\to X+Y$ via one of the four
paths of equation~\eqref{eq:decomp}. Products are one $\mathcal{O}_{1}$
element (gluon) and one $\mathcal{O}_{2}$ element ($W^{\pm}$ or $Z^{0}$),
or three $\mathcal{O}_{1}$ elements.
\item[\textbf{Level 2.}] Secondary decay: $\mathcal{O}_{2}$ elements decompose
into two $\mathcal{O}_{1}$ elements: $(0,1,1)\to(0,1,0)+(0,0,1)$, etc.
\emph{Neutrinos appear at this level.}
\end{enumerate}

\begin{observation}[Neutrino exit at level 2]\label{obs:nu-level}
When an $\mathcal{O}_{2}$ generator decomposes into two $\mathcal{O}_{1}$
elements, one of the products is typically a charged lepton $(Q\neq0)$ and
one is a neutrino ($Q=0$, $R_{\mathrm{surf}}(\nu)\approx0$). The neutrino
exits the coherence network immediately (Observation~\ref{obs:neutrinos}).
The energy it carries -- approximately $12.2\%$ of $M_{H}$ in total -- is
lost to experimental reconstruction. This is not a measurement limitation but
a structural consequence: neutrinos have no surface coherence budget, so they
cannot form further triangles. The fact that the measured Higgs mass is
reproduced at $1.49\sigma$ by equation~\eqref{eq:mh} suggests that the
prediction already accounts, at least approximately, for the mass of the
complete decay state, not just the detectable fraction.
\end{observation}

%%------------------------------------------------------------
\section{Open problems}
%%------------------------------------------------------------

\begin{openproblem}[OP-D62-1: correction on $v$]\label{op:1}
Derive from C1--C4 the correction factor $k_{1}/\Ntot^{2}=9/14161$ on the vev
(Remark~\ref{rem:vcorr}). The candidate interpretation -- weight of the null
endomorphism in $\mathbb{Z}_{\Ntot}$ -- must be formalised as a lemma. Expected
impact: reduction of the vev deviation from $676\ppm$ to $41\ppm$.
\end{openproblem}

\begin{openproblem}[OP-D62-2: derivation of $\lambda_H$]\label{op:2}
Derive the Higgs self-coupling $\lambda_{H}=(M_{H}/v)^{2}/2\approx0.129$
from C1--C4. Prerequisite: top quark mass from OP-D59-1.
\end{openproblem}

\begin{openproblem}[OP-D62-3: self-consistency of the $M_Z$ equation]\label{op:3}
Prove analytically that equation~\eqref{eq:mz} has a unique fixed point, and
that its solution converges from any physically reasonable initial value.
Numerical verification (convergence in two iterations) is established but a
formal proof is missing.
\end{openproblem}

\begin{openproblem}[OP-D62-4: hadronic running coupling from C1--C4]\label{op:4}
Derive $\Dahad=\ksurf/\varphi$ analytically as the hadronic contribution to
$\Delta\alpha(\MZ)$ with lower bound $\mproton$. Then derive the pion and kaon
corrections (from OP-D59-1) to bring $\Dahad^{(\mathrm{PDL})}$ into agreement
with the PDG dispersive value $0.027661\pm0.000110$. Expected impact on $\MZ$:
reduction of the tension from $3.67\sigma$ to below $2\sigma$.
\end{openproblem}

\begin{openproblem}[OP-D62-5: Higgs production cross-section]\label{op:5}
Derive the factor $5^{2}=(\sqrt{5})^{4}$ in equation~\eqref{eq:xsec} from a
combinatorial matching polynomial on $K_4$, connecting the double $K_4$ surface
encounter to the $\Qphi$ structure of the \PDL\ surface leakage parameter.
\end{openproblem}

%%------------------------------------------------------------
\section{Summary and conclusions}
%%------------------------------------------------------------

\begin{table}[h!]
\centering
\caption{Summary of D62 results. Masses in units of $\mproton=938.272\MeV$.
PDG~2024 experimental values. Tensions use experimental uncertainties.}
\label{tab:summary}
\renewcommand{\arraystretch}{1.35}
\begin{tabular}{p{1.5cm}p{5.0cm}ccc}
\toprule
Quantity & \PDL\ expression & PDL value & Exp.\ value & Status\\
\midrule
$\MZ/\MW$ & $1/\cos(19\pi/119)$ & $1.14049$ & $1.13461$ & Theorem, $5182\ppm$\\
$v/\mproton$ & $\Ntot^{2}/(R_{e}k_{1})$ & $262.241$ & $262.418$ & Conj., $676\ppm$\\
$\MZ/\mproton$ & analytical chain & $97.179$ & $97.187$ & Conj., $3.67\sigma$\\
$\MW/\mproton$ & $\MZ\cos(19\pi/119)$ & $85.208$ & $85.656$ & Corollary\\
$\MH/\mproton$ & $\varphi k_{2}/\sinthW+\Delta n\,k_{2}/\Ntot$ & $133.611$ & $133.437$ & Conj., $1.49\sigma$\\
\bottomrule
\end{tabular}
\end{table}

This document makes five contributions to the \PDL\ corpus.

\textbf{Definition of gauge boson mass} (Definition~\ref{def:boson-mass}).
The first formal definition of boson mass in the \PDL\ programme, bridging
the gap between the gauge group derivation (D57--D61) and the mass question.
Gauge bosons are transport modes in $\Ztwo$, not fermionic closures; their
mass is a minimal excitation threshold, not a stationary coherence cost.

\textbf{Derivation of the vev} (Conjecture~\ref{conj:vev}). The electroweak
vev is expressed as a ratio of network capacities: $\vev=\Ntot^{2}/(R_{e}k_{1})
\cdot\mproton$, with each factor justified by an unconditional theorem. The
isolation of this candidate in the numerical search is robust.

\textbf{The $M_Z/M_W$ theorem} (Theorem~\ref{thm:mzmw}). The ratio
$\MZ/\MW=1/\cosW$ is an unconditional theorem of C1--C4. The tree-level PDL
Weinberg angle agrees with the physical MS-bar value to $0.48\sigma$.

\textbf{The $M_Z$ analytical chain} (Conjecture~\ref{conj:mz}). The chain
$\MZ=g_{Z}(\Dalep+\Dahad)\cdot\vev/2$ yields $84.6\ppm$ at $3.67\sigma$.
The tension is precisely attributed: $\Dahad=\ksurf/\varphi$ uses $\mproton$
as the lower integration bound, excluding pion and kaon contributions that
will be recovered when OP-D59-1 is resolved.

\textbf{The decay cascade} (Observations~\ref{obs:neutrinos},
\ref{obs:nu-level}). The Higgs cascade is mapped exactly onto the $\Ztwo$
topology. The four primary decay paths are derived combinatorially. Neutrinos
exit the coherence network at level~2, carrying $12.2\%$ of $M_{H}$ and
rendering the reconstructed LHC mass a coherent-fraction measurement.

The resolution of OP-D59-1 (quark and hadron mass spectrum from C1--C4) is the
key next step: it will supply the missing pion and kaon contributions to
$\Dahad$, derive $\lambda_{H}$ via the top quark Yukawa coupling, and enable
D62's conjectures to be elevated to theorems.

\bibliographystyle{unsrt}
\bibliography{D62}

\end{document}
